\section{How would you know the checker is correct?}
\label{part:correctness}

``Is it correct?'' is two questions wearing one coat, and separating them is most
of the work.

\begin{enumerate}[leftmargin=1.6em]
\item \textbf{Is the specification right?} Does \eqref{eq:ni} capture what you
      mean by confidentiality? This is \emph{not} provable. It is a definition.
\item \textbf{Is the implementation right with respect to the specification?}
      This is provable, in layers, at wildly varying cost.
\end{enumerate}

Conflating them produces the characteristic formal-methods overclaim: a
mechanized proof of question 2 presented as if it settled question 1.

\subsection{Validating the specification: countermodels, not proofs}

You cannot prove a definition. What you can do is show that plausible
\emph{alternative} definitions are wrong, which is what a countermodel is for: a
small executable program witnessing that some appealing principle fails.

The metatheory names seven such required negative results, and they are worth
reading as a validation \emph{method} rather than a list. Each is a principle
somebody would naturally adopt, plus a program refuting it:

\begin{center}
\begin{tabular}{@{}p{0.52\linewidth} p{0.38\linewidth}@{}}
\toprule
Appealing but invalid principle & Refuted by \\
\midrule
bounded-run filtering is a sound proof domain
  & a secret selecting a trip count past the bound \\[2pt]
a future release may condition an earlier observation
  & emit-then-release, two operations \\[2pt]
local per-template checks determine global trace order
  & two events, same payloads, opposite order \\[2pt]
an ordinary SSA slice is closed around a $\phi$
  & the predecessor-choice diamond \\[2pt]
unproved ABI alias separation may be assumed
  & store through $p$, load through $q$, $p=q$ admitted \\[2pt]
an empty premise domain supports a meaningful proof
  & contradictory preconditions \\[2pt]
unary refinement preserves confidentiality
  & a refinement safe alone, unsafe when paired \\
\bottomrule
\end{tabular}
\end{center}

Four of these are encoded in the harness today. The value of encoding them is
that they are the tests a \emph{re-implementation} would fail --- they capture the
reasoning errors, not the coding errors, and reasoning errors are the ones that
survive code review.

\begin{remark}
This is the strongest available answer to ``how do I know the property is the
right property.'' It is an argument from the absence of known counterexamples to
the definition, strengthened by having actively searched for them. That is
weaker than a proof and much stronger than an assertion.
\end{remark}

\subsection{A ladder of assurance for the implementation}

Order these by cost per unit of confidence, cheapest first. The interesting
observation is that the cheapest three are rarely done, and they subsume most of
what people hope to get from the expensive one.

\subsubsection*{Rung 1: metamorphic and differential testing}

Properties of the checker itself, requiring no ground truth:

\begin{itemize}[nosep,leftmargin=1.4em]
\item \textbf{determinism} --- two runs on one input produce byte-identical
      reports. Worth a deliberate worklist-shuffling switch so the guarantee is
      falsifiable rather than accidental; every pointer-keyed hash map on the
      report path is a hazard here.
\item \textbf{permutation invariance} --- renaming SSA values or reordering
      independent operations must not change the verdict. Note this splits in
      two: the \emph{verdict} must be invariant, while the serialized blocker
      sequence must be \emph{canonical}. A single test comparing whole reports
      conflates them.
\item \textbf{monotone refusal} --- obfuscating the input may turn \Verified{}
      into \Unknown{}, never the reverse.
\item \textbf{idempotence} --- re-running on the checker's own normalized output
      agrees.
\end{itemize}

\subsubsection*{Rung 2: an exhaustive small-domain oracle}

This is the highest-value step available and the one most often skipped.

For \emph{small} bit-widths the property is decidable by brute force. Take a
generated program over, say, 4-bit integers and a handful of memory cells.
Enumerate every pair of secret inputs whose entry states are coupled --- equal
public configuration, equal public alias topology, equal initially-visible
components --- run both lanes concretely, and compare the projected traces step by
step, routing release events through the prefix-causal transition of
\cref{sec:ledger}. That yields \textbf{ground truth} --- not an approximation ---
for that program.

\begin{remark}[The oracle must not use the whole-run encoding either]
It is tempting to write the enumeration the short way: keep only the secret pairs
whose end-of-run authorized releases agree, then compare traces. That is
\cref{rem:mtcm3} all over again, and it makes the oracle agree with a buggy
checker rather than catch it --- including on the harness's own
\fx{prefix\_causal\_release.bad.mlir}. An oracle built on the refuted encoding
validates nothing. Couple at entry; retire obligations as releases occur.
\end{remark}

Now you have a soundness oracle, and you can fuzz against it:

\begin{itemize}[nosep,leftmargin=1.4em]
\item generate random programs in the fragment plus random manifests;
\item compute the brute-force verdict and the checker's verdict;
\item any program where brute force finds a leak while the checker reports
      \Verified{} is \textbf{a soundness bug}, reduced automatically to a minimal
      fixture;
\item the reverse direction is not a bug but a precision measurement, and
      tracking its rate over time is the only honest way to report precision.
\end{itemize}

Two things make this unusually effective here. First, the analyzable fragment is
a \emph{closed-world whitelist}, so a generator covering it is a finite,
enumerable engineering task rather than an open-ended one. Second, every
soundness bug it finds arrives as a concrete, minimal, checked-in regression
fixture --- exactly the artifact the corpus wants. A mechanized proof, by
contrast, tells you a theorem holds but hands you no fixture.

\subsubsection*{Rung 2b: cross-level differential}

Rung 2 fuzzes one level against ground truth. The cheaper sibling fuzzes one
level against \emph{another level of the same program}: analyse before and after
each lowering and compare verdicts. No oracle is needed, because the two verdicts
are each other's oracle.

\Verified{} above and \Unsafe{} below is laundering --- the lowering introduced
the leak. The reverse is evidence destruction. Degradation to \Unknown{} is
fine. This is the check that catches the entire class in
Part~\ref{part:lowering}, it reuses the corpus unchanged, and it is the reason
that part argues for extraction at several levels rather than one.

\subsubsection*{Rung 3: certificate-producing analysis}

Rather than proving the analyzer correct, have it emit a \emph{certificate} that
an independent, much smaller checker validates: an unsatisfiability proof from the
solver, the product encoding, and the mapping from IR sites to product
constraints.

The payoff is a shrunken trusted computing base. You stop needing to trust a
large, fast-changing analysis and start needing to trust a small checker plus the
semantics. This is the same trade that makes proof-carrying code and validated
translation attractive, and it has a decisive practical advantage: the analysis
can be rewritten freely without invalidating anything, because each run
re-establishes its own validity.

\subsubsection*{Rung 4: mechanized proof of the core}

Worth doing --- for a carefully chosen core, discussed next.

\subsection{Mechanizing it for MLIR: what is and is not well-posed}

Start with the uncomfortable part. \textbf{``Prove this correct for MLIR'' is not
a well-posed goal.} MLIR is a meta-IR: an extensible framework whose dialects are
open and user-defined. There is no fixed language to have a semantics for, so
there is nothing to quantify over. Any mechanization must first fix a dialect
subset and give \emph{that} a semantics.

The good news is that this project has already done that, without framing it as a
prerequisite for proof. The normal form of \cref{step:normalform} \emph{is} the
fixed subset: a closed-world whitelist of types, opcodes, constants, and
intrinsics. That is the object to mechanize.

\begin{remark}[Reusable prior work, checked]
\label{rem:priorwork}
The foundation is in better shape than the framing above suggests, and it points
at LLVM-IR level rather than MLIR level.

\emph{Vellvm} is a mechanized semantics for a large sequential subset of LLVM IR,
actively maintained against Rocq, and since its 2021 rewrite it is structured
around interaction trees rather than the original 2012 relational semantics.
Decisively for this programme, \emph{noninterference has already been given a
semantics in the same interaction-tree idiom}, by overlapping authors. That means
the two halves --- an LLVM semantics and a noninterference semantics --- exist in
one compatible formalism, which is the single largest reuse opportunity
available and a strong argument for expressing the fragment at LLVM-IR level.

\emph{Alive2} does bounded translation validation of LLVM optimizations via SMT.
Not a proof, but exactly the right shape for validating the normalizer, per
milestone~6 of the ladder below. (``Milestone'' rather than ``step'' throughout
this part: \emph{Step} is reserved for the numbered steps of
Part~\ref{part:walkthrough}.)
\emph{CompCert} proves functional correctness of a C compiler: architecturally
instructive, different property. Its constant-time-preserving successor is the
closer analogue and is discussed in \cref{sec:positioning}.

For MLIR specifically, mechanized semantics exist in Lean and, separately, as
dialects carrying semantics lowered to SMT --- the latter has already found
miscompilations upstream. Neither carries an information-flow theory.
\end{remark}

\subsubsection*{The decomposition that makes it tractable}

Split the goal three ways and give each part the treatment it can actually
support:

\begin{center}
\begin{tabular}{@{}l p{0.42\linewidth} p{0.28\linewidth}@{}}
\toprule
Part & Content & Treatment \\
\midrule
(a) & \eqref{eq:ni} is the right property & definition; reviewed and
      countermodel-validated, \emph{not} proved \\[3pt]
(b) & the product construction implies the property & mechanized once; this is
      the core theorem \\[3pt]
(c) & this run's artifacts are valid & per-run certificate checking, not proof \\
\bottomrule
\end{tabular}
\end{center}

Almost all the leverage is in refusing to prove (c). The analysis, the
normalizer, and the solver integration will change constantly; proving them once
means re-proving them forever. Validating each run costs a fixed engineering
investment and survives arbitrary rewrites.

\subsubsection*{The core theorem, stated}

The load-bearing obligation is:

\begin{theorem}[Product soundness]
\label{thm:core}
For every artifact $\artifact$ in the fragment, every entry $\entry$, and every
coalition $\coal$: if no bad state of the two-lane product is reachable from the
coupled initial states, then $\Sec_\coal$ holds at $\entry$ --- that is,
$\ProdSafe_\coal(\artifact,\entry) \Rightarrow \Sec_{\coal,\entry}(\artifact)$.
\end{theorem}

Everything else the checker does is a strategy for deciding the antecedent. This
direction is already proved: it is Theorem R8.3 of the metatheory, with the full
premise chain assembled in R9.1 and the solver form in R9.2 --- an accepted
\textsc{unsat} for the identity-bound bad-reachability formula establishes
$\Sec$, and \textsc{unknown} does not.

\subsubsection*{The completeness direction is available, and unstated}

It is tempting to say that only soundness is needed and that completeness is
deliberately abandoned in exchange for the right to answer \Unknown{}. That is
not quite the situation, and the difference is worth a theorem.

The metatheory's Theorem R6.3, \emph{reachable-bad correspondence}, is a
\textbf{biconditional}: a bad state is reachable in the product \emph{if and only
if} there exist admitted states and coupled choices, satisfying low-equivalence,
whose unary executions have a first mismatch forbidden by the observation
relation. The $\Leftarrow$ direction is the soundness of
\cref{thm:core}. The $\Rightarrow$ direction, contraposed, is precisely
$\Sec_\coal \Rightarrow \ProdSafe_\coal$ --- the completeness direction --- and
the corpus never states it as a theorem, corollary, or remark.

One qualifier is load-bearing rather than decorative. R6.3 quantifies over a
\emph{first mismatch forbidden by the active prefix-causal relation}, not over
pairs pre-filtered by whole-run release equality. So the corollary below holds
for $\Sec_\coal$ read prefix-causally, as constructed in \cref{sec:ledger}, and
is \textbf{false} under the naive whole-run reading of \eqref{eq:ni}: the
\code{emit-then-release} program satisfies the whole-run biconditional while
leaking. Stating the completeness result without that qualifier would convert a
correct theorem into an unsound one.

\begin{theorem}[Relative completeness, harvestable from R6.3]
\label{thm:complete}
For an artifact in the fragment, with every contract present, within the declared
execution bound, and relative to the fixed observation relation:
$\ProdSafe_\coal \Leftrightarrow \Sec_\coal$.
\end{theorem}

Read that carefully, because the qualifiers carry it. It says the product is
\emph{exact} within the bound: no false negatives, and --- the part a type system
can never claim --- no false positives either. It does \emph{not} say the
observation relation captures every physical channel; the metatheory is explicit
that completeness here is relative to the normal-form grammar and the declared
observation relation, and does not establish that the relation contains every
physical timing, microarchitectural, operating-system, or network observation.

This is the sharpest available statement of what the architecture buys over a
flow-\emph{in}sensitive type system, which carries one fixed level per variable
and therefore rejects programs as simple as ``assign a secret, then overwrite it
with a constant, then publish.'' (A flow-sensitive system, with a per-program-point
environment, accepts that particular program --- it is their motivating example.
What no type system of either kind gets is the converse direction below.)
It costs no new proof --- only the discipline of stating a corollary that already
follows.

\subsubsection*{Where completeness actually leaks}

Every \Unknown{} the checker can return corresponds to a distinct formal reason
the biconditional's antecedent fails to be decidable:

\begin{center}
\begin{tabular}{@{}p{0.46\linewidth} p{0.44\linewidth}@{}}
\toprule
Formal cause & Reason-code family \\
\midrule
undecidability of reachability, so a bound is required
  & bound exhausted, resource limit \\
fragment membership: no semantics outside the whitelist
  & unsupported opcode, type, intrinsic, residual vector \\
missing \emph{relational} contract for a callee
  & missing contract, indirect call, recursion \\
solver incompleteness
  & solver timeout \\
the semantics is a \emph{set} of behaviours
  & freeze may choose, ambiguous NaN payload, uninitialized load \\
memory abstraction cannot decide aliasing
  & alias binding mismatch, layout-dependent pointer comparison \\
channel absent from the observation relation
  & unsupported address observation profile \\
\bottomrule
\end{tabular}
\end{center}

So the reason-code enumeration is not an error taxonomy. It is the exception list
of \cref{thm:complete}, which has a design consequence: the list must be
\emph{closed}, and each member owes an obligation of the form \emph{if this cause
does not apply, the checker decides}.

\begin{remark}[The half that is missing]
\label{rem:falsealarm}
The enumeration covers \emph{refusals}. It does not cover \emph{conservative
rejections} --- cases where the analysis reports a problem on a program that
satisfies the property. At least three such sources exist and none has a code:
over-approximated may-alias, an observation relation strictly stronger than
output-only noninterference, and the requirement of world-level constant
structure even at hosts no coalition can see. Because they carry no code, a false
alarm from any of them is indistinguishable in the report from a genuine finding.
If the completeness bookkeeping is to be as disciplined as the soundness
bookkeeping, these need first-class representation --- the same argument that
produced the precision controls of \cref{sec:pairs}, arriving from the theory
side rather than the testing side.
\end{remark}

A second, smaller theorem is worth mechanizing because it is where the
architecture could quietly go wrong:

\begin{theorem}[The diagnostic layer only narrows]
\label{thm:diag}
No diagnostic disposition implies $\Sec_\coal$, and no diagnostic disposition
removes a product constraint. In particular, a relational-required disposition
implies nothing whatever.
\end{theorem}

\cref{thm:diag} is small, and it is exactly the invariant that
\cref{rem:danger} says implementations break under pressure to reduce false
positives.

\subsubsection*{A milestone ladder}

\begin{enumerate}[leftmargin=1.8em]
\item \textbf{Mechanize the fragment's operational semantics}, with the
      observation projection as an explicit output rather than an afterthought.
      Deliverable: a step relation and a trace projection. The projection being
      explicit is what makes the property statable at all.
\item \textbf{State \eqref{eq:ni} and the product formally}, and prove
      \cref{thm:core}. This is the single highest-value proof in the programme,
      and it is independent of any implementation.
\item \textbf{Prove \cref{thm:diag}} for the actual disposition domain.
\item \textbf{Extract or validate the product encoder.} Either extract a verified
      encoder from the proof, or keep the fast one and check its output against
      the mechanized semantics per run. Prefer the latter.
\item \textbf{Do not mechanize the solver.} Consume unsatisfiability certificates
      from it and check those. Verifying an SMT solver is a research programme in
      its own right and buys little here.
\item \textbf{Validate the normalizer per run}, translation-validation style:
      show the normalized module is observationally equivalent to its input on
      this input, rather than proving the normalizer once. It will change
      constantly; the proof would not survive.
\item \textbf{Leave the backend to bounded validation.} Comparing the emitted
      machine code's observable trace structure against the IR-level one, on the
      actual artifact, is the only thing that addresses the
      trace-preservation obligation of \cref{step:deployment}. A verified backend
      would be better and is a far larger undertaking.
\end{enumerate}

\subsubsection*{What a proof can never give you}

Every \Lv{4} obligation is outside the reach of any proof about the IR: real
helper latency, whether the backend preserved the trace, hardware isolation
between tenants, speculative execution, sanitizer sufficiency, certificate
soundness. A mechanized theorem is always \emph{relative to a declared machine
model}, and the gap between that model and the silicon is exactly what the
\Conditional{} outcome exists to record.

This is worth stating loudly because it is the most common way formal claims
mislead. ``Proved noninterferent'' with an unstated machine model is not a
stronger claim than ``\Conditional{}, pending these four named facts'' --- it is
the same claim with the assumptions hidden. The four-outcome contract is, in that
light, less a limitation of this system than an unusually honest interface.

\subsection{Where this sits in the literature}
\label{sec:positioning}

Worth knowing precisely, because two things here are already owned and two are
not.

\paragraph{The product construction is classical.} Self-composition as a route to
noninterference is Barthe, D'Argenio and Rezk (CSFW 2004, extended in
\emph{MSCS} 2011); the framing of secure information flow as a safety problem is
Terauchi and Aiken (SAS 2005); the general theory of properties over sets of
traces is Clarkson and Schneider's hyperproperties (CSF 2008, \emph{JCS} 2010);
and relational Hoare logic is Benton (POPL 2004). For the type-system side of the
two-layer split, the canonical soundness result is Volpano, Irvine and Smith
(\emph{JCS} 4(2--3):167--187, 1996), surveyed by Sabelfeld and Myers
(\emph{IEEE JSAC} 2003).

\paragraph{Constant-time verification at LLVM level is occupied.} \emph{ct-verif}
(Almeida, Barbosa, Barthe, Dupressoir, Emmi; USENIX Security 2016) is the direct
ancestor: product construction, LLVM, annotations for public inputs. Two facts
about it matter. Its toolchain is effectively abandoned --- last commit 2017 ---
so it cannot be run as a baseline. But its benchmark corpus survives and is still
the field's de-facto suite: twelve directories covering \texttt{curve25519-donna},
\texttt{openssl}, \texttt{mee-cbc}, \texttt{polarssl}, \texttt{sodium},
\texttt{s2n} and others. \textbf{Reuse the corpus as fixtures; do not expect to
reuse the tool.} The live successor is \emph{CT-LLVM} (Zhang and Barthe, 2025),
which is analysis-based --- def-use plus alias analysis --- rather than
product-based, so it makes a different soundness/precision trade.

\paragraph{Compiler preservation is occupied.} The mechanized analogue of the
deployment obligation in \cref{step:deployment} is CompCert-CT, ``Formal
verification of a constant-time preserving C compiler'' (Barthe, Blazy,
Grégoire, Hutin, Laporte, Pichardie, Trieu; POPL 2020). The optimization-level
statement is Rosemann, Hack and Garg, ``Non-interference Preserving Optimising
Compilation'' (OOPSLA 2025), which develops hyperproperty simulations --- close
kin to the paired-refinement notion this system uses. The general framework for
what such a theorem should even say is Abate et al., ``Journey Beyond Full
Abstraction'' (CSF 2019).

\paragraph{What is not occupied.} Two gaps, and the second is larger.

First, a \emph{mechanized, product-construction} noninterference checker at LLVM
IR with a preservation story. The pieces exist separately --- Vellvm for the
semantics, the interaction-tree noninterference work for the property, CompCert-CT
for preservation --- and nobody has assembled them.

Second, and more striking: \textbf{MLIR has no information-flow story at all.}
There is no security, information-flow, taint, secrecy, or constant-time dialect
upstream. The nearest neighbour by name, HEIR's \texttt{secret} dialect, is a
lifting mechanism for homomorphic encryption; its own documentation states that
\texttt{secret.generic} does not itself do anything, and it carries no
noninterference property, no declassification, and no theorem. Mechanized MLIR
semantics exist, and information flow has not been built on them.

That is the opening, and it argues for a specific division of labour: prove the
core at LLVM-IR level where the semantics already exists, and treat MLIR as the
place where the \emph{policy surface} and the tensor-level channels live.

\subsection{If you only do three things}

\begin{enumerate}[leftmargin=1.8em]
\item Build the \textbf{exhaustive small-domain oracle} and fuzz against it. It
      finds real soundness bugs now, it produces regression fixtures as output,
      and it requires no proof assistant.
\item Make the checker \textbf{emit a certificate} and write the small
      independent validator. This shrinks what you have to trust, and it survives
      every future rewrite of the analysis.
\item Mechanize \textbf{\cref{thm:core} only}. One theorem, no implementation
      dependencies, and it is the statement everything else is a strategy for.
\end{enumerate}

Rungs 1--3 are weeks to months. A mechanized \cref{thm:core} over a small
fragment is months. A verified end-to-end pipeline is years, and would still
terminate in the same \Lv{4} obligations you started with.
